\documentclass[../../main.tex]{subfiles}% 注意这里的文件路径不能用 ./main.tex ,否则用latexmk编译子文件会报错
\graphicspath{{\subfix{./image/}}{\subfix{../../image/}}} % 指定图片引用路径,后续可以直接使用图片文件名
% 如果图片存放在上级目录,则使用 ../../image/ 作为备用路径

% 例如:
% \begin{figure}[H]
% \centering
% \includegraphics[width=0.3\textwidth]{图.png}
% \caption{}
% \label{figure:图}
% \end{figure}
% 注意:上述\label{}一定要放在\caption{}之后,否则引用图片序号会只会显示??.

\begin{document}
% SOURCE p=75; printed=60; object=ch02-convergence-modes-title; region=bottom
\section{收敛方式}\label{section:ch02-modes-of-convergence}
% SOURCE p=75,76; printed=60,61; object=ch02-convergence-modes-intro; region=bottom/top
若 $\{f_n\}$ 为集合 $X$ 上的复值函数序列，陈述“当 $n\to\infty$ 时，$f_n\to f$”可以有许多不同含义，例如逐点收敛或一致收敛。若 $X$ 为测度空间，还可以讨论几乎处处收敛或 $L^1$ 收敛。当然，一致收敛蕴含逐点收敛，后者又蕴含几乎处处收敛（一般不能反过来）；但是，这些收敛方式不蕴含 $L^1$ 收敛，反之也不成立。记住下面这些在 $\mathbb{R}$ 上（配备 Lebesgue 测度）的例子很有用。
% SOURCE p=76; printed=61; object=ch02-convergence-examples; region=upper
\begin{examples}\label{examples:ch02-convergence-modes}
\begin{enumerate}[label=\roman*.,ref=(\roman*)]
\item\label{item:ch02-convergence-spreading} $f_n=n^{-1}\chi_{(0,n)}$。
\item\label{item:ch02-convergence-translating} $f_n=\chi_{(n,n+1)}$。
\item\label{item:ch02-convergence-concentrating} $f_n=n\chi_{[0,\frac1n]}$。
\item\label{item:ch02-convergence-typewriter} $f_1=\chi_{[0,1]}$，$f_2=\chi_{[0,\frac12]}$，$f_3=\chi_{[\frac12,1]}$，$f_4=\chi_{[0,\frac14]}$，$f_5=\chi_{[\frac14,\frac12]}$，$f_6=\chi_{[\frac12,\frac34]}$，$f_7=\chi_{[\frac34,1]}$；一般地，$f_n=\chi_{[\frac{j}{2^k},\frac{j+1}{2^k}]}$，其中 $n=2^k+j$，$0\le j<2^k$。
\end{enumerate}
在\subcref{examples:ch02-convergence-modes}{item:ch02-convergence-spreading}、\subcref{examples:ch02-convergence-modes}{item:ch02-convergence-translating}和\subcref{examples:ch02-convergence-modes}{item:ch02-convergence-concentrating}中，$f_n\to0$ 分别为一致收敛、逐点收敛和几乎处处收敛，但 $f_n$ 不在 $L^1$ 中收敛到 $0$（事实上，对所有 $n$，$\int|f_n|=\int f_n=1$）。在\subcref{examples:ch02-convergence-modes}{item:ch02-convergence-typewriter}中，$f_n\to0$ 在 $L^1$ 中成立，因为当 $2^k\le n<2^{k+1}$ 时，$\int|f_n|=2^{-k}$；但对任意 $x\in[0,1]$，$f_n(x)$ 都不收敛，因为使 $f_n(x)=0$ 的 $n$ 有无穷多个，使 $f_n(x)=1$ 的 $n$ 也有无穷多个。
\end{examples}
% SOURCE p=76; printed=61; object=ch02-dominated-lone-convergence; region=middle
\begin{proposition}\label{proposition:ch02-dominated-lone-convergence}
若 $f_n\to f$ 几乎处处成立，且对所有 $n$，$|f_n|\le g\in L^1$，则 $f_n\to f$ 在 $L^1$ 中成立。
\end{proposition}
% SOURCE p=76; printed=61; object=ch02-proof-dominated-lone-convergence; region=middle
\begin{proof}
由于 $|f_n-f|\le2g$，这由\cref{theorem:ch02-dominated-convergence}显然得出。
\end{proof}
% SOURCE p=76; printed=61; object=ch02-subsequence-intro; region=middle
另外，下面将看到，若 $f_n\to f$ 在 $L^1$ 中成立，则存在子序列几乎处处收敛于 $f$。另一种常用的收敛方式是依测度收敛。
% SOURCE p=76; printed=61; object=ch02-cauchy-in-measure; region=middle
\begin{definition}[依测度 Cauchy]\label{definition:ch02-cauchy-in-measure}
设 $\{f_n\}$ 为 $(X,\mathscr{M},\mu)$ 上的可测复值函数序列。若对每个 $\epsilon>0$，
\begin{align*}
\mu\bigl(\{x:|f_n(x)-f_m(x)|\ge\epsilon\}\bigr)\longrightarrow0\quad(m,n\to\infty),
\end{align*}
则称 $\{f_n\}$ \textbf{依测度 Cauchy}。
\end{definition}
% SOURCE p=76; printed=61; object=ch02-convergence-in-measure; region=middle-lower
\begin{definition}[依测度收敛]\label{definition:ch02-convergence-in-measure}
若对每个 $\epsilon>0$，
\begin{align*}
\mu\bigl(\{x:|f_n(x)-f(x)|\ge\epsilon\}\bigr)\longrightarrow0\quad(n\to\infty),
\end{align*}
则称 $\{f_n\}$ \textbf{依测度收敛}于 $f$。
\end{definition}
% SOURCE p=76; printed=61; object=ch02-measure-convergence-examples; region=middle-lower
\begin{remark}
例如，上述\subcref{examples:ch02-convergence-modes}{item:ch02-convergence-spreading}、\subcref{examples:ch02-convergence-modes}{item:ch02-convergence-concentrating}和\subcref{examples:ch02-convergence-modes}{item:ch02-convergence-typewriter}中的序列都依测度收敛于零，但\subcref{examples:ch02-convergence-modes}{item:ch02-convergence-translating}不是依测度 Cauchy 的。
\end{remark}
% SOURCE p=76; printed=61; object=ch02-lone-implies-measure; number=2.29; region=lower
\begin{proposition}\label{proposition:ch02-lone-implies-measure}
若 $f_n\to f$ 在 $L^1$ 中成立，则 $f_n\to f$ 依测度成立。
\end{proposition}
% SOURCE p=76; printed=61; object=ch02-proof-lone-implies-measure; region=lower
\begin{proof}
令 $E_{n,\epsilon}=\{x:|f_n(x)-f(x)|\ge\epsilon\}$。则 $\int|f_n-f|\ge\int_{E_{n,\epsilon}}|f_n-f|\ge\epsilon\mu(E_{n,\epsilon})$，所以 $\mu(E_{n,\epsilon})\le\epsilon^{-1}\int|f_n-f|\to0$。
\end{proof}
% SOURCE p=76; printed=61; object=ch02-measure-not-lone; region=lower
\begin{remark}
\cref{proposition:ch02-lone-implies-measure}的逆命题不成立，\subcref{examples:ch02-convergence-modes}{item:ch02-convergence-spreading}和\subcref{examples:ch02-convergence-modes}{item:ch02-convergence-concentrating}即为反例。
\end{remark}
% SOURCE p=76; printed=61; object=ch02-measure-cauchy-convergence; number=2.30; region=lower
\begin{theorem}\label{theorem:ch02-measure-cauchy-convergence}
设 $\{f_n\}$ 依测度 Cauchy。则存在可测函数 $f$，使 $f_n\to f$ 依测度成立，且存在子序列 $\{f_{n_j}\}$ 几乎处处收敛于 $f$。此外，若还满足 $f_n\to g$ 依测度成立，则 $g=f$ 几乎处处成立。
\end{theorem}
% SOURCE p=76,77; printed=61,62; object=ch02-proof-measure-cauchy-convergence; region=bottom/top
\begin{proof}
可选取 $\{f_n\}$ 的子序列 $\{g_j\}=\{f_{n_j}\}$，使得若 $E_j=\{x:|g_j(x)-g_{j+1}(x)|\ge2^{-j}\}$，则 $\mu(E_j)\le2^{-j}$。若 $F_k=\bigcup_{j=k}^\infty E_j$，则 $\mu(F_k)\le\sum_k^\infty2^{-j}=2^{1-k}$；且若 $x\notin F_k$，则对 $i\ge j\ge k$，有
% SOURCE p=76; printed=61; object=ch02-measure-cauchy-estimate; number=2.31; region=bottom
\begin{align}\label{align:ch02-measure-cauchy-estimate}
|g_i(x)-g_j(x)|\le\sum_{l=j}^{i-1}|g_{l+1}(x)-g_l(x)|\le\sum_{l=j}^{i-1}2^{-l}\le2^{1-j}.
\end{align}
% SOURCE p=77; printed=62; object=ch02-proof-measure-cauchy-convergence-continued; region=top
故 $\{g_j\}$ 在 $F_k^c$ 上逐点 Cauchy。令 $F=\bigcap_1^\infty F_k=\limsup E_j$。则 $\mu(F)=0$；若在 $x\notin F$ 时令 $f(x)=\lim g_j(x)$，在 $x\in F$ 时令 $f(x)=0$，则 $f$ 可测（见\cref{exercise:ch02-convergence-set,exercise:ch02-measurable-pasting}），且 $g_j\to f$ 几乎处处成立。此外，\eqref{align:ch02-measure-cauchy-estimate}表明，当 $x\notin F_k$ 且 $j\ge k$ 时，$|g_j(x)-f(x)|\le2^{1-j}$。由于 $k\to\infty$ 时 $\mu(F_k)\to0$，所以 $g_j\to f$ 依测度成立。但这也蕴含 $f_n\to f$ 依测度成立，因为
\begin{align*}
\{x:|f_n(x)-f(x)|\ge\epsilon\}
&\subseteq\{x:|f_n(x)-g_j(x)|\ge\tfrac12\epsilon\}\\
&\qquad\cup\{x:|g_j(x)-f(x)|\ge\tfrac12\epsilon\},
\end{align*}
且当 $n$ 和 $j$ 充分大时，右端两个集合的测度都很小。同样，若 $f_n\to g$ 依测度成立，则对所有 $n$，
\begin{align*}
\{x:|f(x)-g(x)|\ge\epsilon\}
&\subseteq\{x:|f(x)-f_n(x)|\ge\tfrac12\epsilon\}\\
&\qquad\cup\{x:|f_n(x)-g(x)|\ge\tfrac12\epsilon\},
\end{align*}
所以对所有 $\epsilon$，$\mu(\{x:|f(x)-g(x)|\ge\epsilon\})=0$。让 $\epsilon$ 沿某个数列趋于零，即得 $f=g$ 几乎处处成立。
\end{proof}
% SOURCE p=77; printed=62; object=ch02-lone-ae-subsequence; number=2.32; region=middle
\begin{corollary}\label{corollary:ch02-lone-ae-subsequence}
若 $f_n\to f$ 在 $L^1$ 中成立，则存在子序列 $\{f_{n_j}\}$，使 $f_{n_j}\to f$ 几乎处处成立。
\end{corollary}
% SOURCE p=77; printed=62; object=ch02-proof-lone-ae-subsequence; region=middle
\begin{proof}
结合\cref{proposition:ch02-lone-implies-measure}和\cref{theorem:ch02-measure-cauchy-convergence}。
\end{proof}
% SOURCE p=77; printed=62; object=ch02-egoroff-intro; region=middle
若 $f_n\to f$ 几乎处处成立，并不能推出 $f_n\to f$ 依测度成立，\subcref{examples:ch02-convergence-modes}{item:ch02-convergence-translating}说明了这一点。不过，在有限测度空间上，这一结论确实成立，而且还有更强的结果。
% SOURCE p=77; printed=62; object=ch02-egoroff; number=2.33; region=middle-lower
\begin{theorem}[Egoroff 定理]\label{theorem:ch02-egoroff}
设 $\mu(X)<\infty$，$f_1,f_2,\ldots$ 和 $f$ 为 $X$ 上的可测复值函数，且 $f_n\to f$ 几乎处处成立。则对每个 $\epsilon>0$，存在 $E\subseteq X$，使 $\mu(E)<\epsilon$，且 $f_n\to f$ 在 $E^c$ 上一致收敛。
\end{theorem}
% SOURCE p=77; printed=62; object=ch02-proof-egoroff; region=lower
\begin{proof}
不妨假设 $f_n\to f$ 在 $X$ 上处处成立。对 $k,n\in\mathbb{N}$，令
\begin{align*}
E_n(k)=\bigcup_{m=n}^\infty\{x:|f_m(x)-f(x)|\ge k^{-1}\}.
\end{align*}
对固定的 $k$，$E_n(k)$ 随 $n$ 递增而递减，且 $\bigcap_{n=1}^\infty E_n(k)=\varnothing$；由于 $\mu(X)<\infty$，可得当 $n\to\infty$ 时 $\mu(E_n(k))\to0$。给定 $\epsilon>0$ 及 $k\in\mathbb{N}$，选取充分大的 $n_k$，使 $\mu(E_{n_k}(k))<\epsilon2^{-k}$，并令 $E=\bigcup_{k=1}^\infty E_{n_k}(k)$。则 $\mu(E)<\epsilon$，且当 $n>n_k$、$x\notin E$ 时，有 $|f_n(x)-f(x)|<k^{-1}$。因此 $f_n\to f$ 在 $E^c$ 上一致收敛。
\end{proof}
% SOURCE p=77; printed=62; object=ch02-almost-uniform; region=bottom
\begin{definition}[几乎一致收敛]\label{definition:ch02-almost-uniform}
Egoroff 定理结论中出现的收敛方式，有时称为\textbf{几乎一致收敛}。
\end{definition}
% SOURCE p=77; printed=62; object=ch02-almost-uniform-implications; region=bottom
\begin{proposition}\label{proposition:ch02-almost-uniform-implications}
不难看出，几乎一致收敛蕴含几乎处处收敛和依测度收敛（\cref{exercise:ch02-almost-uniform-implications}）。
\end{proposition}
\subsection*{习题}
% SOURCE p=78; printed=63; object=ch02-exercise-measure-metric; number=32; region=top
\begin{exercise}\label{exercise:ch02-measure-metric}
设 $\mu(X)<\infty$。若 $f$ 和 $g$ 为 $X$ 上的复值可测函数，定义
\begin{align*}
\rho(f,g)=\int\frac{|f-g|}{1+|f-g|}\,\mathrm{d}\mu.
\end{align*}
若将几乎处处相等的函数视为相同，则 $\rho$ 是可测函数空间上的一个度量；且 $f_n\to f$ 关于这个度量收敛，当且仅当 $f_n\to f$ 依测度成立。
\end{exercise}
% SOURCE p=78; printed=63; object=ch02-exercise-fatou-in-measure; number=33; region=upper
\begin{exercise}\label{exercise:ch02-fatou-in-measure}
若 $f_n\ge0$，且 $f_n\to f$ 依测度成立，则 $\int f\le\liminf\int f_n$。
\end{exercise}
% SOURCE p=78; printed=63; object=ch02-exercise-dominated-in-measure; number=34; region=upper
\begin{exercise}\label{exercise:ch02-dominated-in-measure}
设 $|f_n|\le g\in L^1$，且 $f_n\to f$ 依测度成立。
\begin{enumerate}
\item\label{item:ch02-dominated-measure-integrals} $\int f=\lim\int f_n$。
\item\label{item:ch02-dominated-measure-lone} $f_n\to f$ 在 $L^1$ 中成立。
\end{enumerate}
\end{exercise}
% SOURCE p=78; printed=63; object=ch02-exercise-measure-one-epsilon; number=35; region=upper-middle
\begin{exercise}\label{exercise:ch02-measure-one-epsilon}
$f_n\to f$ 依测度成立，当且仅当对每个 $\epsilon>0$，存在 $N\in\mathbb{N}$，使得在 $n\ge N$ 时，$\mu(\{x:|f_n(x)-f(x)|\ge\epsilon\})<\epsilon$。
\end{exercise}
% SOURCE p=78; printed=63; object=ch02-exercise-characteristic-lone-limit; number=36; region=upper-middle
\begin{exercise}\label{exercise:ch02-characteristic-lone-limit}
若对 $n\in\mathbb{N}$，$\mu(E_n)<\infty$，且 $\chi_{E_n}\to f$ 在 $L^1$ 中成立，则 $f$（几乎处处等于）某个可测集的示性函数。
\end{exercise}
% SOURCE p=78; printed=63; object=ch02-exercise-composition-convergence; number=37; region=middle
\begin{exercise}\label{exercise:ch02-composition-convergence}
设 $f_n$ 和 $f$ 为可测复值函数，且 $\phi:\mathbb{C}\to\mathbb{C}$。
\begin{enumerate}
\item\label{item:ch02-composition-ae} 若 $\phi$ 连续，且 $f_n\to f$ 几乎处处成立，则 $\phi\circ f_n\to\phi\circ f$ 几乎处处成立。
\item\label{item:ch02-composition-uniform} 若 $\phi$ 一致连续，且 $f_n\to f$ 一致收敛、几乎一致收敛或依测度收敛，则 $\phi\circ f_n\to\phi\circ f$ 分别一致收敛、几乎一致收敛或依测度收敛。
\item\label{item:ch02-composition-counterexamples} 当 $\phi$ 不满足相应连续性假设时，存在反例。
\end{enumerate}
\end{exercise}
% SOURCE p=78; printed=63; object=ch02-exercise-measure-algebra; number=38; region=middle-lower
\begin{exercise}\label{exercise:ch02-measure-algebra}
设 $f_n\to f$ 和 $g_n\to g$ 都依测度成立。
\begin{enumerate}
\item\label{item:ch02-measure-sums} $f_n+g_n\to f+g$ 依测度成立。
\item\label{item:ch02-measure-products} 若 $\mu(X)<\infty$，则 $f_ng_n\to fg$ 依测度成立；但若 $\mu(X)=\infty$，则未必如此。
\end{enumerate}
\end{exercise}
% SOURCE p=78; printed=63; object=ch02-exercise-almost-uniform-implications; number=39; region=lower
\begin{exercise}\label{exercise:ch02-almost-uniform-implications}
若 $f_n\to f$ 几乎一致收敛，则 $f_n\to f$ 几乎处处收敛且依测度收敛。
\end{exercise}
% SOURCE p=78; printed=63; object=ch02-exercise-egoroff-dominated; number=40; region=lower
\begin{exercise}\label{exercise:ch02-egoroff-dominated}
在\cref{theorem:ch02-egoroff}中，可以将假设“$\mu(X)<\infty$”替换为“对所有 $n$，$|f_n|\le g$，其中 $g\in L^1(\mu)$”。
\end{exercise}
% SOURCE p=78; printed=63; object=ch02-exercise-sigma-finite-uniform-pieces; number=41; region=lower
\begin{exercise}\label{exercise:ch02-sigma-finite-uniform-pieces}
若 $\mu$ 为 $\sigma$-有限测度，且 $f_n\to f$ 几乎处处成立，则存在可测集 $E_1,E_2,\ldots\subseteq X$，使得 $\mu((\bigcup_1^\infty E_j)^c)=0$，且 $f_n\to f$ 在每个 $E_j$ 上一致收敛。
\end{exercise}
% SOURCE p=78; printed=63; object=ch02-exercise-counting-uniform; number=42; region=lower
\begin{exercise}\label{exercise:ch02-counting-uniform}
设 $\mu$ 为 $\mathbb{N}$ 上的计数测度。则 $f_n\to f$ 依测度成立，当且仅当 $f_n\to f$ 一致收敛。
\end{exercise}
% SOURCE p=78; printed=63; object=ch02-exercise-parameter-almost-uniform; number=43; region=bottom
\begin{exercise}\label{exercise:ch02-parameter-almost-uniform}
设 $\mu(X)<\infty$，且函数 $f:X\times[0,1]\to\mathbb{C}$ 满足：对每个 $y\in[0,1]$，$f(\cdot,y)$ 可测；对每个 $x\in X$，$f(x,\cdot)$ 连续。
\begin{enumerate}
\item\label{item:ch02-parameter-tail-measurable} 若 $0<\epsilon,\delta<1$，则 $E_{\epsilon,\delta}=\{x:|f(x,y)-f(x,0)|\le\epsilon\text{ 对所有 }y<\delta\text{ 成立}\}$ 可测。
\item\label{item:ch02-parameter-tail-uniform} 对任意 $\epsilon>0$，存在集合 $E\subseteq X$，使 $\mu(E)<\epsilon$，且当 $y\to0$ 时，$f(\cdot,y)\to f(\cdot,0)$ 在 $E^c$ 上一致收敛。
\end{enumerate}
\end{exercise}
% SOURCE p=79; printed=64; object=ch02-exercise-lusin; number=44; region=top
\begin{exercise}[Lusin 定理]\label{exercise:ch02-lusin}
若 $f:[a,b]\to\mathbb{C}$ 为 Lebesgue 可测函数，且 $\epsilon>0$，则存在紧集 $E\subseteq[a,b]$，使 $\mu(E^c)<\epsilon$，且 $f|_E$ 连续。（利用\cref{theorem:ch02-egoroff}和\cref{theorem:ch02-lone-density}。）
\end{exercise}
\end{document}
